\section{Experimental Protocol}
\label{sec:experiments}

\subsection{Data Governance}
Table~\ref{tab:data} gives the complete data inventory at the selected 3-s
horizon. Splits are by recording, not by overlapping windows. In the fixed
split, hyperparameters and blend weights use only the two real validation
recordings; the real test recording is excluded from gradient fitting and the
initial fixed-split architecture selection. Multiple post-selection
diagnostics (cross-validation ensembles, stillness gating and blends) were
subsequently evaluated on the same recording, so these comparisons are
exploratory rather than a strict one-shot confirmatory test. In the
leave-one-session-out protocol, every one of the eight non-test recordings is
held out in turn; the four rigidity recordings that passed the mount test
(rotation-dominant motion with still phases, 2723 pairs) are training-only
in every fold and never scored. Synthetic seeds are disjoint across train,
validation, test, and OOD.

\begin{table}[t]
\centering
\caption{Datasets at the selected 3-s horizon with 2-s context.}
\label{tab:data}
\footnotesize
\setlength{\tabcolsep}{4pt}
\begin{tabular}{lrrrl}
\toprule
Dataset & Seq. & Duration & Pairs & Role \\
\midrule
Real train & 6 & 495.5~s & 6754 & Fit (fixed split) \\
Real validation & 2 & 260.8~s & 1436 & Selection (fixed split) \\
Real train$+$validation & 8 & 756.3~s & 8190 & 8 CV folds \\
Real extra-train (rigid) & 4 & 342~s & 2723 & Training only \\
Real test & 1 & 115.2~s & 647 & held-out test \\
Synthetic train & 60 & 1200~s & 7200$^\ast$ & Pretrain / direct \\
Synthetic validation & 15 & 300~s & 1800$^\ast$ & Twin audit \\
Synthetic test & 15 & 300~s & 2400 & Test \\
Synthetic OOD & 5 & 100~s & 800 & Stress test \\
\bottomrule
\multicolumn{5}{l}{$^\ast$Pairs capped per sequence to balance motion families.}
\end{tabular}
\end{table}

\subsection{Baselines, Measurements, and Metrics}
The nonlearned baselines are zero motion, zero-initial-velocity double
integration, and ridge regression on acceleration/gyroscope summary features.
We adapt official-source architectures from RoNIN (ResNet and
LSTM)~\cite{yan2020ronin}, TLIO (ResNet mean head)~\cite{liu2020tlio}, and
IMUNet (depthwise-separable CNN)~\cite{zeinali2024imunet}. Their input/output
layers are changed from pedestrian 2-D velocity to the common six-axis, 3-D,
3-s displacement protocol. These results are denoted \emph{adapted/retrained};
they are not official pretrained-model scores. Official repositories contain
no directly compatible weights, and TLIO explicitly requires retraining on a
user-generated dataset. Our earlier Conv--BiGRU estimator on demeaned
body-frame channels is kept as an ablation of the representation.

Every architecture is evaluated under the same three regimes: synthetic
direct, real-only, and synthetic-pretrained/real-fine-tuned. IMUNet, the
strongest public architecture on this protocol, is additionally trained
under the leave-one-session-out protocol with the same extra-train data as
PhysNet (two seeds per fold, sixteen models each).
Parameter counts include all trainable parameters. Latency is averaged per
window after warm-up on an NVIDIA GeForce RTX 5090 Laptop GPU with batch size
16 and CUDA synchronization; it is a relative architecture comparison, not an
embedded-device measurement.

\paragraph{Estimator-independent measurements.}
Two experiments characterize the protocol rather than a model and use only
reference data and the raw inertial stream. The information-limit measurement
computes, for every pair, the mean camera velocity over the context from the
reference positions, and compares $\bar{\bf v}\Delta t$ with the residual of
the estimator and with its own magnitude as a function of context length. The
zero-velocity bound detects still phases in the reference (speed below
3~mm/s for at least 0.4~s), integrates the gravity-compensated acceleration
between consecutive phases with the velocity pinned at both ends, solves the
one constant bias per interval that this pinning identifies, and compares the
result with the reference displacement; the attitude is taken either from the
reference (upper bound) or from the integrated gyroscope (operational case).
Both are reported per interval length so that the crossover between
model-based and learned estimation is visible.

\paragraph{Metrics.}
Local displacement error is
\begin{equation}
 e_p=\frac{1}{N}\sum_{n=1}^{N}
 \|\widehat{\Delta{\bf p}}_n-\Delta{\bf p}_n\|_2 ,
\end{equation}
where $N$ is the number of displacement pairs. We also report the
multivariate coefficient of determination
$R^2=1-\sum_n\|\widehat{\Delta{\bf p}}_n-\Delta{\bf p}_n\|_2^2/
\sum_n\|\Delta{\bf p}_n-\bar{\Delta{\bf p}}\|_2^2$, where
$\bar{\Delta{\bf p}}=N^{-1}\sum_n\Delta{\bf p}_n$; the error in the fastest
third, defined by the top third of reference displacement magnitudes; and the median
predicted magnitude (a shrinkage indicator). Because overlapping windows are
strongly correlated, confidence intervals on the held-out test use a block
bootstrap over non-overlapping 3-s blocks (2000 resamples). Trajectory error
uses anchor-only absolute trajectory error
$\mathrm{ATE}=[M^{-1}\sum_{i=1}^{M}\|\hat{\bf p}_i-{\bf p}_i\|_2^2]^{1/2}$
over the $M$ finite pose-graph nodes; no post-hoc rigid alignment is applied
because each connected component already has one metric anchor. For
the digital twin, the audit reports ideal/noisy kinematic consistency,
timestamp statistics, corner detection rate, corner-to-exact-projection RMSE,
and split-seed overlap.

\paragraph{Ablations.}
All design choices are made on the fixed validation split only: for the
earlier Conv--BiGRU, context and horizon sweeps, loss weighting, pooling,
camera-aligned axes, twin version and replay ratio; for PhysNet, with two
seeds each, the variants listed in Table~\ref{tab:ablation}.
Section~\ref{sec:results} reports every post-selection evaluation performed on
the held-out recording.
